% Supplementary material. Compiles standalone; same style resolution as main.tex.
\documentclass[10pt,twocolumn,letterpaper]{article}
\IfFileExists{cvpr.sty}{\usepackage{cvpr}}{\usepackage{cvpr_fallback}}
\usepackage{amsmath,amssymb}
\usepackage{graphicx}
\usepackage{booktabs}
\usepackage[colorlinks,linkcolor=blue!60!black,citecolor=blue!60!black]{hyperref}
\newcommand{\BandUseHigh}{9}
\newcommand{\BandUseLow}{45}
\newcommand{\Ceiling}{41.9}
\newcommand{\FloorHigh}{0.072}
\newcommand{\FloorLow}{0.036}
\newcommand{\GmacAtBudget}{338}
\newcommand{\IntraGmac}{453.5}
\newcommand{\MainHighRate}{18.8}
\newcommand{\MainLowRate}{33.2}
\newcommand{\MainMean}{25.4}
\newcommand{\MapBits}{89}
\newcommand{\MapOverheadLow}{0.011}
\newcommand{\NumSeq}{40}
\newcommand{\SatHigh}{0.396}
\newcommand{\SatLow}{0.179}
\newcommand{\SatWindowHi}{0.194}
\newcommand{\SatWindowMb}{15}
\newcommand{\SeamArlsHigh}{0.197}
\newcommand{\SeamLinearHigh}{0.384}
\newcommand{\SeamReplHigh}{0.218}
\newcommand{\SeamZerosHigh}{0.677}
\newcommand{\TilingOverhead}{8.5}
\newcommand{\WallMasked}{9.9}
\newcommand{\WallPredicted}{35.3}
\newcommand{\WallSorted}{28.5}

\newcommand{\dB}{\,dB}

\title{Where to Stop: Tile-Adaptive Early Exit in a Learned Image Decoder\\
       \large Supplementary Material}
\author{Anonymous CVPR submission\\ Paper ID ****}

\begin{document}
\maketitle

\section{Architecture details}

\paragraph{Costs.} With $K$ exits over $N{=}12$ trunk blocks and split depth $j$,
the cost of exit $k$ in units of one released decode is
\begin{equation}
c_k = s_{\mathrm{up}} + s_{\mathrm{trunk}}\frac{(k{+}1)b}{N}
      + s_{\mathrm{head}} + a_k + r,
\end{equation}
where $b = N/K$ blocks per exit, $s_{\mathrm{up}} = 0.0816$,
$s_{\mathrm{trunk}} = 0.8944$, $s_{\mathrm{head}} = 0.0240$ are the measured
shares of the released decoder, $a_k$ is the adapter's share and $r$ the
seam-repair module's $0.0095$. Only the trunk term depends on $k$: the stem and
head run for every tile at every depth, which is why the ceiling
$100(1-c_j)$ is well below $100$ however shallow the ladder gets.

\paragraph{Adapters.} $\mathrm{Ad}(f) = f + Wf$ with $W\in\mathbb{R}^{C\times C}$
zero-initialised costs $C^2$ MAC/px, one eighth of a block's $8C^2 + 9C$. The
FFN variant, $\mathrm{Ad}(f)=f+\mathrm{PW}(\mathrm{WSiLUChunkAdd}
(\mathrm{PW}_{C\to 4C}(f)))$, costs $2C^2$. Assignment is by the number of blocks
the exit skips: FFN at four or more, the $1\times1$ otherwise.

\section{The tiling penalty in detail}

\paragraph{Why only the depthwise matters.} A \texttt{DepthConvBlock} at $C=384$
costs $8C^2 + 9C = 1{,}183{,}104$ MAC/px, of which the $3\times3$ depthwise is
$9C = 3{,}456$ --- $0.29\%$. Every other operator is $1\times1$ and has no
neighbours to miss. This is why the exact remedy is affordable: haloing or
coupling only the depthwise costs $+0.032\%$ of the decode at $256$\,px tiles,
against $1.745\times$ for a halo carried through the whole block, which an
earlier iteration of this work measured and rejected --- correctly about the
number and wrongly about the target.

\paragraph{Contamination is areal, not linear.} With $b$ per-tile blocks on a
tile of side $F$ feature pixels the corrupted fraction is
$1 - ((F-2b)/F)^2$. At our $F=32$, $b=8$ this is $0.750$. The small-$b$
approximation $4b/F$ gives $1.000$ here and should not be used.

\paragraph{The repair module's ceiling.} Splitting per-pixel squared error by
distance from the nearest tile boundary at $q63$, with repair on against off and
the same exit map:

\begin{center}
\begin{tabular}{lrrr}
\toprule
Band (px) & Share & OFF & Change \\
\midrule
0--4    & 6.2\%  & $4.7\times10^{-5}$ & $-0.27\%$ \\
4--16   & 17.3\% & $4.4\times10^{-5}$ & $+0.05\%$ \\
16--64  & 51.6\% & $4.3\times10^{-5}$ & $+0.04\%$ \\
64--128 & 25.0\% & $4.1\times10^{-5}$ & $+0.04\%$ \\
\bottomrule
\end{tabular}
\end{center}

The error is $15\%$ higher on the boundary band than in the interior and falls
monotonically with distance, so the artefact is real and local. But the gate
leaks: the module gains only in the first band. Taking the boundary gain and
assuming a perfect gate elsewhere bounds what it could ever earn at
$0.062 \times 4.7\!\times\!10^{-5} \times 0.0027 / 4.30\!\times\!10^{-5}
\approx 1.8\times10^{-4}$ of the frame's error, i.e. $\approx 0.0008\dB$, for
$0.95\%$ of the decode.

\section{Reproducibility}

\paragraph{Frozen encoder.} Every evaluation asserts
$\max|\theta^{\mathrm{enc}}_{\mathrm{ours}} - \theta^{\mathrm{enc}}_{\mathrm{rel}}| = 0$
before measuring. Without it the comparison is not on the same latent and the
claim of an unchanged bitstream is unverified.

\paragraph{Warm-start equivalence.} At initialisation the deepest exit
reproduces the released decoder bit-exactly at $q0/32/63$, because every
inherited weight is copied and every new module's last layer is zero.

\paragraph{Sorted execution.} The reordering of
Section~5 is bit-identical on CUDA under a random exit map and random content.
It is \emph{not} bit-identical on CPU, and the reason is not the algorithm: a
plain \texttt{DepthConvBlock} is already order-dependent at batch size 13 there,
by about $10^{-7}$, while exact at 4, 9, 16 and 40, because oneDNN selects a
different blocking. A test that compared the two paths on a decode of an
all-zero image would pass vacuously --- every tile carries the same latent, so
any permutation is trivially exact --- and ours therefore uses random content.

\paragraph{Two decibel conventions.} Pooling every tile of every frame into one
MSE and averaging a per-frame decibel differ by $0.023$--$0.033\dB$ on an
identical allocation, a quarter to a third of the working budget, with pooling
always the flattering one. We report per-frame, which is what the reference
implementation of the baseline computes.

\paragraph{Two quality metrics.} An RGB MSE ratio and the $6{:}1{:}1$-weighted
YUV 4:2:0 PSNR the baseline reports agree within $0.003\dB$ on identical
decodes across the whole rate range; we verified this rather than assuming it.

\section{Negative and withdrawn results}

We list these because they cost real time and are not otherwise recoverable from
the paper.

\begin{itemize}\itemsep2pt
\item \textbf{First-order border extrapolation.} Continuing the local gradient
      past a tile boundary is $1.8\times$ worse than a zeroth-order hold at
      $q63$ (\SeamLinearHigh$\dB$ against \SeamReplHigh). Amplifying boundary
      noise costs more than the trend recovers.
\item \textbf{Fitted AR(1) padding}~\cite{arls}. The best estimator we measured,
      and rejected: $0.02\dB$ for $10.7\%$ of decode wall-clock.
\item \textbf{Learned seam repair.} Bounded above at $\approx 0.0008\dB$ for
      $0.95\%$ of decode, even with a perfect gate.
\item \textbf{Jointly training the router with the decoder.} In one
      configuration the router collapsed to a rate-dependent constant with zero
      agreement with the oracle; in another it did not. Training the router
      against a frozen decoder is reliable and is what we report.
\item \textbf{A halo through the per-tile trunk.} $1.745\times$ the cost of a
      full decode --- more than routing saves.
\end{itemize}

{\small
\bibliographystyle{ieee_fullname}
\bibliography{refs}
}
\end{document}
